\Needspace{10\baselineskip}
\section{新增证据留下的具体可研究问题}
\begin{longtable}{P{36mm}P{67mm}P{46mm}}
\toprule 问题 & 现在证据到哪里 & 合适的下一项计算\\\midrule\endfirsthead
\toprule 问题 & 现在证据到哪里 & 合适的下一项计算\\\midrule\endhead
全局率是否趋于平台 & 固定宏观形状时率随尺寸下降，较大Hall域下降变缓；平均区间存在稿件矛盾\citep{F600} & 固定片厚与固定纵横比两套扫描，分别报告局地/全局率、通量阶段和重复扰动误差\\
Hall与PIC何时足够相近 & 强导引场合并中宏观率接近，压力分解和次级岛仍不同\citep{L241} & 用十矩及非局地闭合扫描导引场、离子温度和尺寸，不在每例重新调闭合参数\\
磁岛何时控制净能量转换 & 多种有岛/无岛结果可有相近率；数值耗散与质量比会改变岛\citep{F606,F788} & 同步收敛拓扑、率、压力应变及边界通量；比较形成时间与被排走时间\\
非对称电子层怎样局域化 & 不同非理想项导致不同层长与岛序列\citep{F625} & 匹配电子尺度厚度，再比较张量闭合与超电阻，并单独扫描域长\\
电子单独重联的二次生成 & 大尺度喷流与电子剪切两条路径有二维时序支持\citep{F655} & 多格/$d_e$与滤波收敛、鞍点轨迹、三维及电子温度自由度\\
Full-F爆发与暂态的关系 & 非正规性与非线性谱展宽均出现，但因果量化不足\citep{L314} & 用物理自由能范数算最优增益，控制初始扰动并测非线性转折\\
Hall发电机符号由何控制 & 多螺旋后续事件出现竞争，样本有限\citep{F798} & 保留实际压力剖面和壁面预算，扩大事件数并按前态分类\\
多组分与有限装置 & 冷离子改变分离线电场，中性拖曳降速而层宽未扩\citep{L227,L190} & 多离子/中性独立动量模型，匹配实际边界和碰撞尺度\\
\bottomrule
\end{longtable}

这些问题中，速率、层长、压力输运与大尺度耦合可用Hall或扩展流体继续推进；电子非热尾、微观不可逆加热、尘埃充电和完整多束分布则需额外方程或动力学对照。模型能生成某种形态，不等于其所有微观解释已包含在模型里。

\section{原文一致性审计：保留疑点，不替作者猜参数}
\label{sec:remaining-audit}
本表针对实际核读版本。它记录需要读者注意的公式或标注问题，不表示已证明作者程序错误；版本不同也可能已修正。报告推导均从明确符号和基本方程重建。
\begin{longtable}{P{24mm}P{76mm}P{49mm}}
\toprule 文献/位置 & 原文问题 & 本报告处理\\\midrule\endfirsthead
\toprule 文献/位置 & 原文问题 & 本报告处理\\\midrule\endhead
F600 v1页3--4 & 图1平均磁通阴影约0.5--3，正文写0.05--0.3 & 引用尺寸趋势，精确平均统计待澄清\\
F589 v1页7 & 张角写为$L_i/d_m$，与前文定义及随后公式的倒数关系不符 & 使用$d_m/L_i$，明确为几何重建\\
F606 v1式4 & 若$e>0$，电子压力与惯性项符号应与电子动量相符，印刷式有疑点 & 用本章电子欧姆定律统一符号\\
L314 v2页6 & 写$k=1$同时给$\Delta'=14.3108$，与$2(1/k-k)$不一致 & 不将该组文字参数作为可复现输入\\
L314 v2图11--12 & 坐标写$\hat d_e^2$，图注写$\hat d_e$且行列描述反置 & 保留参数映射不确定性；不据此精确比较跨排序区域\\
E473 v1图6 & 正文与图注左右对应相反 & 依据图内标签：左Full-F，右Full-F+LWL\\
E473/L314能量式 & 磁能写在种类求和中可能重复；个别系数/符号不一致 & 从正定极化二次型及物理磁能重建，不直接抄表\\
F655 v1页6、附录A & 电子/离子Alfvén速度比写反；主事件时刻320.7与360.7混用 & 采用$V_{Ae}/V_{Ai}=\sqrt{m_i/m_e}$；具体数据配对待核验\\
L283 v2页2 & $d_{i0}$式呈现为频率量纲，扰动坐标与$xz$几何不一致 & 用$c/\omega_{pi}$定义尺度，不擅自补齐复现初值\\
F703 v1正文/附录A & 干涉零点基准视线的文字对应不一致 & 保留诊断系统误差，不把投影唯一反演为三维Hall结构\\
\bottomrule
\end{longtable}

正式版F788对噪声、有限电流片寿命和全局几何的讨论比早期稿更谨慎，故本报告以正式版为主要依据。L314引言对F788“欠分辨”的方向性转述与原研究不符，本报告采用F788自己的算例和讨论。这个交叉核对说明，综述不能仅沿引用链复制结论。
